\section{Evaluation}\label{sec:eval}

\paragraph{Setup.}  Apple M4~Pro (8 performance + 4 efficiency cores, 12 hardware
threads, 24~GB unified memory), macOS~15, Apple~Clang \texttt{-O3}, LLVM~OpenMP.
Baseline is UPPAAL~5.0.0 \texttt{verifyta} with default options (breadth-first,
conservative space optimisation, minimal constraint systems), invoked as
\texttt{verifyta -s -q}.  Timings are wall-clock, measured externally to the
process so parsing and setup are charged to every tool alike.  Peak memory is the
kernel's maximum resident set (\texttt{getrusage}, via \texttt{/usr/bin/time -l}),
the same metric for all three engines.  \tool{} \emph{omp} uses 12 threads;
\emph{metal} is the GPU backend on the same chip.

\paragraph{Benchmarks.}  We use scalable families from the standard UPPAAL and
LTSmin benchmark sets, at sizes large enough for the measurement to be
meaningful.  We report stored-state counts as a correctness check: \tool{} and
UPPAAL agree exactly on every row of \cref{tab:main}, which is the strongest
available evidence that the two perform the same exploration.  We exclude
families whose query is discharged after a handful of states (\emph{csma},
\emph{hddi}, \emph{single}, \emph{train} all terminate in under
\SI{50}{\milli\second} for every tool and measure nothing), and \emph{grc} and
\emph{twolock}, where all three engines exceed our timeout at every size we
tried---they are hard for UPPAAL and for us alike, and we have no result to
report.

\begin{table}[t]
\centering
\caption{Wall-clock time (seconds) and speedup over UPPAAL~5.0.  \emph{omp} is
the 12-thread CPU backend, \emph{metal} the GPU backend.  The \emph{states}
column is the count stored by \emph{omp} and by UPPAAL, which agree exactly on
every row; \emph{metal} agrees on all rows except \emph{phil10}, where its merge
order differs and it stores \num{94776} (see \cref{sec:eval}).  Best engine per
row in bold.}
\label{tab:main}
\setlength{\tabcolsep}{5pt}
\begin{tabular}{lrrrrrrr}
\toprule
& & \multicolumn{3}{c}{time (\si{\second})} & \multicolumn{3}{c}{speedup} \\
\cmidrule(lr){3-5}\cmidrule(lr){6-8}
model & states & UPPAAL & omp & metal
      & $\frac{\text{UPPAAL}}{\text{omp}}$
      & $\frac{\text{UPPAAL}}{\text{metal}}$
      & $\frac{\text{omp}}{\text{metal}}$ \\
\midrule
arbiter15 & \num{278528}  & 2.40  & 1.06          & \textbf{0.68} & $2.3\times$ & $3.5\times$ & $1.6\times$ \\
arbiter16 & \num{589824}  & 4.71  & \textbf{1.87} & 2.09          & $2.5\times$ & $2.3\times$ & $0.9\times$ \\
arbiter17 & \num{1245184} & 11.52 & 4.32          & \textbf{4.31} & $2.7\times$ & $2.7\times$ & $1.0\times$ \\
\midrule
fischer10 & \num{447598}  & 2.82  & 1.23          & \textbf{0.55} & $2.3\times$ & $5.1\times$ & $2.2\times$ \\
fischer11 & \num{1464971} & 11.02 & 5.02          & \textbf{1.94} & $2.2\times$ & $5.7\times$ & $2.6\times$ \\
\midrule
viking12  & \num{241662}  & 0.63  & \textbf{0.14} & 0.31          & $4.5\times$ & $2.0\times$ & $0.5\times$ \\
viking15  & \num{2768894} & 10.11 & \textbf{1.63} & 3.13          & $6.2\times$ & $3.2\times$ & $0.5\times$ \\
\midrule
phil10    & \num{71640}   & 0.53  & \textbf{0.11} & 0.20          & $4.8\times$ & $2.7\times$ & $0.6\times$ \\
\bottomrule
\end{tabular}
\end{table}

\begin{table}[t]
\centering
\caption{Peak resident set (MiB).  The OpenMP backend is memory-competitive with
UPPAAL; the GPU backend is not, by an order of magnitude.}
\label{tab:mem}
\setlength{\tabcolsep}{6pt}
\begin{tabular}{lrrrrr}
\toprule
& \multicolumn{3}{c}{peak RSS (\si{\mebi\byte})} & \multicolumn{2}{c}{ratio} \\
\cmidrule(lr){2-4}\cmidrule(lr){5-6}
model & UPPAAL & omp & metal
      & $\frac{\text{omp}}{\text{UPPAAL}}$
      & $\frac{\text{metal}}{\text{omp}}$ \\
\midrule
arbiter15 & 64  & 72  & \num{1902} & $1.1\times$ & $26\times$ \\
arbiter16 & 106 & 128 & \num{5266} & $1.2\times$ & $41\times$ \\
arbiter17 & 212 & 251 & \num{7071} & $1.2\times$ & $28\times$ \\
fischer10 & 72  & 101 & 920        & $1.4\times$ & $9\times$  \\
fischer11 & 186 & 242 & \num{3822} & $1.3\times$ & $16\times$ \\
viking12  & 53  & 45  & 313        & $0.8\times$ & $7\times$  \\
viking15  & 423 & 219 & \num{3387} & $0.5\times$ & $15\times$ \\
phil10    & 32  & 41  & 210        & $1.3\times$ & $5\times$  \\
\bottomrule
\end{tabular}
\end{table}

\begin{table}[t]
\centering
\caption{Ablation on the CPU backend, 12 threads, wall-clock seconds.
\emph{base} is the engine of \cref{sec:engine} with neither optimisation;
\emph{+gc} adds amortised reclamation (\cref{sec:gc}); \emph{+filter}
additionally enables compressed-form subsumption (\cref{sec:filter}).
Stored-state counts are identical across all three columns on every row except
\emph{phil}, which is nondeterministic for an unrelated reason (see text).}
\label{tab:ablation}
\begin{tabular}{lrrrr}
\toprule
model & base & +gc & +filter & total \\
\midrule
arbiter15 & 1.21 & 1.03 & 0.66 & $1.8\times$ \\
arbiter17 & 9.58 & 6.74 & 3.94 & $2.4\times$ \\
fischer10 & 1.38 & 1.20 & 1.26 & $1.1\times$ \\
fischer11 & 5.95 & 5.08 & 5.10 & $1.2\times$ \\
viking15  & 1.81 & 1.83 & 1.53 & $1.2\times$ \\
phil10    & 0.12 & 0.12 & 0.09 & $1.3\times$ \\
\bottomrule
\end{tabular}
\end{table}


\paragraph{CPU versus GPU versus UPPAAL.}

Both \tool{} backends beat UPPAAL on every model in \cref{tab:main}, by
$2.2$--$6.2\times$ (\emph{omp}) and $2.0$--$5.7\times$ (\emph{metal}).  The two
backends do \emph{not} dominate one another: the GPU is $2.6\times$ faster than
twelve CPU threads on \emph{fischer11} and $1.6\times$ on \emph{arbiter15}, but
roughly \emph{half} the speed on \emph{viking15} and \emph{phil10}.

The discriminating variable is \emph{not} host-side overhead, which is the
intuitive explanation and the wrong one.  Profiling the GPU backend per phase:

\begin{center}
\begin{tabular}{lrrrrr}
\toprule
model & levels & work items & items/dispatch & GPU expand & host share \\
\midrule
fischer11 & 26  & \num{11840631} & \num{455409} & \SI{693}{\milli\second}  & $59\%$ \\
arbiter17 & 19  & \num{12255232} & \num{645012} & \SI{1723}{\milli\second} & $39\%$ \\
viking15  & 100 & \num{8355840}  & \num{83558}  & \SI{1407}{\milli\second} & $43\%$ \\
\bottomrule
\end{tabular}
\end{center}

\emph{fischer11} carries the \emph{largest} host share of the three and still wins
by $2.6\times$, so per-level host cost cannot be what decides the outcome.  Two
effects are visible instead, and they are independent.

First, \textbf{dispatch occupancy}.  The same kernel on the same chip retires
\num{17085} work items per millisecond on fischer11 but only \num{5939} on
viking15 --- a $2.9\times$ difference in per-item throughput, tracking the $5.4\times$
difference in items per dispatch.  With frontiers half as wide spread over five
times as many levels, viking's launches are too small to fill the machine.  The
decisive figure is that viking's GPU expansion \emph{alone} costs
\SI{1.41}{\second} against \SI{1.66}{\second} for the entire CPU run: even with
host and merge time driven to zero the GPU would barely win, so no amount of
host-side tuning changes the verdict there.

Second, \textbf{total GPU-path overhead}, which is what decides arbiter17.  Its
expansion is the best-occupied of the three, yet host and merge add
\SI{2.97}{\second} on top and the CPU backend wins outright.  Occupancy alone
therefore does not predict the winner either: fischer and arbiter sit at
comparable items per dispatch and land on opposite sides.

We draw a deliberately negative conclusion from this.  We attempted a static
routing heuristic (send a model to the GPU when items per dispatch exceeds a
threshold) and abandoned it: any threshold separating viking from fischer also
sends arbiter17 to the slower backend.  Choosing a backend per model requires a
cost model spanning expansion \emph{and} the merge/host path, or an online race
between them; a single occupancy statistic is not sufficient.  We report the two
backends separately rather than their envelope for this reason.

We also measured the CPU-expand/GPU-merge hybrid the unified-memory architecture
permits, in case viking's problem was purely expansion: it improves viking15 from
\SI{3.12}{\second} to \SI{2.60}{\second} but still trails the \SI{1.66}{\second}
CPU backend, and it \emph{regresses} fischer11
($2.61 \to \SI{2.99}{\second}$) and arbiter17 ($4.69 \to \SI{5.35}{\second}$).

\emph{One caveat on state counts.}  On \emph{phil10} the GPU backend stores
\num{94776} symbolic states where the CPU backend and UPPAAL store
${\approx}\num{71600}$.  Both verdicts agree.  The cause is the same
order-sensitivity of convex-union merging discussed in \cref{sec:correctness}:
the GPU merge visits a bucket's candidates in a different order than the CPU
antichain fold, so it performs fewer exact-convex merges and retains more zones.
It is deterministic per backend (three GPU runs gave \num{94776} every time) and
does not occur on the other seven models, where the two backends agree
byte-for-byte.

\paragraph{Memory.}  \Cref{tab:mem} is, we think, the more interesting table.
The CPU backend now sits within $1.4\times$ of UPPAAL's peak RSS on every model
and \emph{below} it on \emph{viking12} and \emph{viking15}
(\SI{219}{\mebi\byte} against \SI{423}{\mebi\byte})---this is the minimal
constraint systems of \cref{sec:mingraph} doing their job, now that
\cref{sec:filter} has made them cheap to query.  Both optimisations of
\cref{sec:opt} leave peak RSS unchanged to within $1\%$
(\SI{249}{\mebi\byte} for arbiter17 before and after), which confirms that
\cref{prop:gcmem} is far from tight in practice and---more tellingly---that the
per-level sweep it replaced was reclaiming very little for its cost.

The GPU backend is another matter: $5$--$41\times$ the CPU backend's memory, and
\SI{7071}{\mebi\byte} on arbiter17 against \SI{251}{\mebi\byte}.  Instrumenting
it attributes the peak to two sources, neither of which is the compression
question this paper is about.  Its passed set is stored \emph{dense} on the host
(\SI{0.27}{\giga\byte} on arbiter15, exactly $n^2$ words per zone), because the
merge kernels consume dense matrices; and its unified-memory pools are sized to
the widest BFS level rather than to a bounded working set
(\SI{1.26}{\giga\byte} on the same model, the single largest component).
Bounding the frontier batch shrinks those pools by $3\times$ as measured at level
boundaries, but leaves the kernel-reported maximum RSS unmoved, so a transient
allocation we have not yet isolated dominates the high-water mark.  We report
this as an open item rather than a result: the GPU backend trades an order of
magnitude of memory for a speedup it only sometimes achieves, and porting
\cref{sec:mingraph}'s representation to it is future work
(\cref{sec:conclusion}).

\paragraph{The ablation is not uniform, and that is the point.}

\Cref{tab:ablation} deserves a careful reading rather than a headline number.

\emph{Amortised reclamation} (\cref{sec:gc}) helps where a run has many BFS
levels over a large intern table---arbiter17 ($9.58 \to \SI{6.74}{\second}$) and
fischer11 ($5.95 \to \SI{5.08}{\second}$)---and does nothing on \emph{viking15}
($1.81 \to \SI{1.83}{\second}$), whose intern table never grows enough for the
per-level sweep to dominate.  It fixes an Amdahl term, so its value tracks how
large that term was.

\emph{Compressed-form subsumption} (\cref{sec:filter}) is dramatic on arbiter17
($6.74 \to \SI{3.94}{\second}$, a further $1.7\times$), useful on viking and
phil, and \emph{neutral on fischer} ($5.08 \to \SI{5.10}{\second}$; and
$1.20 \to \SI{1.26}{\second}$ at fischer10, a slight regression within
run-to-run noise).  The explanation is the one \cref{tab:addprof} predicts: the
filter pays in proportion to how many insertions are duplicates rejected by the
first subsumption test.  Arbiter sits at $89.8\%$; fischer's exploration is
dominated by successor generation and by candidates that survive the first test
and must be unpacked anyway, so the extra $O(|G|)$ scan buys little and costs a
little.  We report this rather than selecting the favourable family: the
optimisation removes a specific cost, and where that cost is absent it correctly
does nothing.

\paragraph{An observation on determinism.}  While collecting
\cref{tab:ablation} we found that \emph{phil10} reports a different stored-state
count on nearly every run: five consecutive 12-thread runs of one binary gave
\numlist{71403;71175;71721;71210;71493}, while three single-threaded runs gave
\num{71470} every time.  This is the order-sensitivity of convex-union merging
discussed in \cref{sec:correctness}, not an artefact of either optimisation---it
predates both, and the verdict is invariant.  We record it because a reader
reproducing \cref{tab:ablation} will otherwise think something is broken, and
because it is a reminder that ``same number of stored states'' is a check
available only on models where the merge order happens not to matter.

\paragraph{Threats to validity.}  (i)~Single machine, single microarchitecture;
the M4~Pro's heterogeneous cores mean the achievable 12-thread ceiling is nearer
$9\times$ than $12\times$, so parallel efficiency computed against $12$
understates the engine.  (ii)~We compare against UPPAAL's default configuration
and have not swept its search-order or state-representation options.  (iii)~Our
state-count agreement with UPPAAL is evidence of equivalent exploration, not a
proof of equivalent semantics; the property we prove (\cref{thm:samestates}) is
internal, relating our engine with and without the fast path.  (iv)~Two families
are excluded for timing out on all engines, so \cref{tab:main} is not a claim
about coverage of the benchmark set.
